{"latex_body":"This section addresses the practical question of how compositional digital systems behave when gates, wires, or state updates are no longer ideal. In the noiseless setting, closure is usually discussed as an algebraic property: a class of functions is stable under composition, substitution, or wiring. Under noise, one asks for a stronger notion of \emph{robust closure}: not only should the intended computation remain in the class, but its implementation should remain reliable under perturbation, faults, and finite redundancy.\n\nA convenient summary statistic is the reliability function\n\\[\n\\varepsilon(r) \\coloneqq 1 - \\Pr[\\text{error}],\n\\]\nas a function of a resource parameter \\(r\\), such as depth, size, redundancy level, code distance, or number of correction rounds. In an idealized analysis, \\(\\varepsilon(r)\\) should improve with \\(r\\) for a well-designed fault-tolerant construction; in practice, it often exhibits a tradeoff between overhead and residual error. The central design problem is therefore to choose a composition scheme whose reliability improves faster than the noise accumulates.\n\n\\paragraph{Noisy gates, majority logic, and redundancy.}\nA basic model assumes that each gate independently fails with some small probability \\(p\\). Even when the underlying logical operation is simple, repeated composition can amplify error unless the architecture includes redundancy. The classical response is to replace a single logical operation by a gadget that computes the same function on multiple copies and then reconciles the outputs by a robust decision rule. Majority logic is the canonical example: if three copies of a bit are available and at most one is corrupted, the majority vote recovers the intended value. More generally, majority-based schemes implement a compositional symmetry: the logical value is invariant under permutations of the replicated wires, and the decision rule factors through this symmetry.\n\nThis symmetry perspective is useful because it explains why redundancy is not merely duplication. The replicated physical states are organized into equivalence classes, and the correction rule is designed to be insensitive to the admissible perturbations within each class. In that sense, majority logic is a small but concrete instance of robust closure: the logical operation survives because the implementation has been quotiented by a noise-tolerant symmetry.\n\n\\paragraph{Von Neumann schemes.}\nVon Neumann's approach to reliable computation formalizes this idea at scale. Instead of trusting a single noisy gate, one builds a circuit from noisy components together with redundancy and periodic correction. The key insight is that reliability can be made to improve recursively if the noise rate is below a threshold and if the redundancy overhead is controlled. In such schemes, the architecture itself becomes part of the computation: the logical function is not merely evaluated, but evaluated through a fault-tolerant encoding and decoding process.\n\nA useful way to read these constructions is as iterative refinement. Each layer takes a noisy physical realization, extracts a more stable logical state, and then re-encodes it into a form that is again suitable for further computation. The resulting design is compositional in the strong sense that the correction mechanism composes with the computation rather than sitting outside it.\n\n\\paragraph{Error-correcting codes as compositional symmetries.}\nError-correcting codes provide a particularly clean language for robust closure. An encoding map embeds logical data into a larger physical space, while a decoding map projects back after noise. The code space is characterized by invariances and syndromes: many distinct physical states represent the same logical state, and the correction procedure identifies which perturbations are admissible. From the perspective of compositional algebra, a code acts like a symmetry-respecting section of a larger system. The logical operations are those that preserve the code space up to correctable error, so the code defines a closure condition relative to the noise model.\n\nThis viewpoint is useful because it separates three layers:\n\\begin{itemize}\n  \\item the logical specification of the computation,\n  \\item the physical implementation subject to noise,\n  \\item the correction and decoding layer that restores the logical semantics.\n\\end{itemize}\nA robust design is one in which these layers compose cleanly, so that the error introduced by one stage remains within the correction capability of the next. In practice, this means that the code distance, the gadget design, and the scheduling of correction rounds must be chosen together rather than independently.\n\n\\paragraph{Threshold theorems for reliable computation.}\nThreshold theorems formalize the possibility of scalable fault tolerance. Informally, they state that if the physical error rate is below a certain threshold, then arbitrarily long computations can be performed reliably with only polylogarithmic or polynomial overhead, depending on the model and assumptions. The theorem is not a claim that noise disappears; rather, it asserts that the logical error can be driven down faster than the circuit grows, provided the architecture is sufficiently redundant and the noise is sufficiently weak.\n\nIn the language of this book, a threshold theorem is a closure result under perturbation: the class of computable functions remains stable when the implementation is lifted from ideal gates to noisy gadgets, as long as the noise lies in a regime where the correction scheme closes the error dynamics. The threshold is therefore a boundary between two qualitative behaviors: below it, recursive correction dominates; above it, error accumulation overwhelms the computation.\n\nA schematic way to express this is to view the logical error rate after one protected layer as a function \\(f(p)\\) of the physical error rate \\(p\\). A threshold regime is one in which \\(f(p) < p\\) for sufficiently small \\(p\\), so repeated concatenation or repeated correction drives the effective error downward. The exact form of \\(f\\) depends on the code family and noise model, but the qualitative fixed-point picture is common across many fault-tolerant constructions.\n\n\\paragraph{Reliability as a design metric.}\nThe function \\(\\varepsilon(r)\\) is useful not only as a theoretical abstraction but also as an engineering metric. For a given family of implementations, one can measure how reliability changes with redundancy, depth, or code distance. Typical plots compare error probability against overhead, showing whether additional resources buy meaningful robustness. Such plots are especially informative when different fault-tolerance strategies are compared on the same benchmark.\n\nIn a well-behaved family, increasing redundancy should improve reliability up to the point where the overhead itself becomes the dominant cost. This makes \\(\\varepsilon(r)\\) a multi-objective quantity: one wants high reliability, but also acceptable size, latency, and energy use. The point of robust closure is not to eliminate tradeoffs, but to ensure that the tradeoffs remain controlled as the system scales.\n\n\\paragraph{Artifact: fault-injection harness.}\nA practical way to study these effects is to build a fault-injection harness that perturbs gates, wires, or state transitions according to a chosen noise model and then measures the resulting output error rate. The harness can sweep over redundancy parameters and produce error-versus-redundancy plots, making the tradeoff between cost and reliability explicit. Such an artifact is valuable for validating analytic claims, for comparing majority-based and code-based schemes, and for identifying regimes where a proposed construction fails to achieve robust closure.\n\nA useful harness reports not only the final error probability but also intermediate statistics such as syndrome rates, correction success rates, and the distribution of failure modes. Those measurements help distinguish between designs that fail because the noise model is too strong and designs that fail because the correction logic is structurally mismatched to the computation.\n\n\\paragraph{Summary.}\nReliability under noise is not an afterthought to compositional design; it is a structural constraint on what counts as a stable implementation. Majority logic, von Neumann-style redundancy, and error-correcting codes all realize the same broad principle: replicate, detect, and correct so that the logical computation survives physical perturbation. Threshold theorems then provide the asymptotic guarantee that this principle can scale. In that sense, robust closure is the noisy analogue of algebraic closure: the system remains within its intended computational class even after the environment has acted on it.","completeness_percent":97,"completeness_rationale":"The draft already covers the imported source themes in full: the reliability function \\(\\varepsilon(r)\\), noisy gates, majority logic, von Neumann schemes, error-correcting codes as compositional symmetries, threshold theorems, and the fault-injection harness artifact. It is coherent with the surrounding book’s emphasis on closure, symmetry, and artifacts, and it preserves the existing prose while refining it into a publishable section draft. The remaining work is mainly citation registration and, if desired, a concrete figure or theorem cross-reference once the artifact set is finalized."}
